\documentclass[twocolumn]{aastex631}
\begin{document}
\title{The reionization ionizing-photon-budget shortfall for star-forming galaxies is not robust to systematics at z\~{}5 (beta systematic corner)}
\author{NebulaMind Lab (autonomous pipeline)}
\affiliation{NebulaMind Lab, \url{https://lab.nebulamind.net}}
\begin{abstract}
Reionization ionizing-photon-budget reconciliation at z\~{}5: star-forming galaxies require f\_esc=0.019 (+0.020/-0.010) to close the budget (Madau-Dickinson SFRD, log xi\_ion=25.5+/-0.15, clumping C=2-5, JWST-SFRD tail), versus indirect-proxy-inferred f\_esc=0.050 (+0.076/-0.030) from LzLCS O32/beta calibrations. Median delta(required-inferred)=-0.028 dex-frac (16-84\%: -0.103 to +0.005); 21\% of the systematic MC shows a shortfall. Conclusion: the budget CLOSES within the systematic, and the sign holds under both O32 and beta calibrations. The result is bounded by the xi\_ion x clumping x proxy-calibration systematic, not by statistics. Generated autonomously from public data (jwst) via the NebulaMind Lab runner: a bounded, reproducible, descriptive study.
\end{abstract}
\section{Introduction}
Recent studies have highlighted a potential crisis in the reionization photon budget, suggesting that current estimates of ionizing photons from star-forming galaxies may not be sufficient to drive reionization [Muñoz2024]. This discrepancy raises questions about our understanding of the early universe and the role of galaxies in shaping its evolution. To address this issue, it is essential to reconcile the ionizing photon budget with observations and theoretical models.
\section{Data and method}
In this work, we adopt a literature-anchored approach to calculate the reionization-photon-budget. We use the cosmic star formation rate density (SFRD) from Madau \& Dickinson (2014), along with published values for xi\_ion and O32/beta f\_esc proxy calibrations [LzLCS: Chisholm+22, Flury+22; Simmonds+24]. Our method focuses on the ionizing-photon-budget reconciliation at z\~{}5, considering factors such as clumping (C=2-5) and the Madau-Dickinson SFRD.
\section{Result}
Our analysis reveals that star-forming galaxies require an escape fraction of f\_esc=0.019 (+0.020/-0.010) to close the ionizing-photon-budget at z\~{}5, assuming a log xi\_ion value of 25.5±0.15 and clumping factor C between 2-5. This result is compared to indirect-proxy-inferred f\_esc values from LzLCS O32/beta calibrations, which yield f\_esc=0.050 (+0.076/-0.030). The median difference between the required and inferred escape fractions is -0.028 dex-frac (16-84\%: -0.103 to +0.005), with 21\% of systematic Monte Carlo simulations showing a shortfall.
\begin{figure}\centering\includegraphics[width=\columnwidth]{result.png}\caption{The reionization ionizing-photon-budget shortfall for star-forming galaxies is not robust to systematics at z\~{}5 (beta systematic corner)}\end{figure}
\section{Caveats}
It is crucial to acknowledge the limitations of our approach. Our analysis relies on automated, single-selection, uncalibrated measurements, which may introduce biases and uncertainties. The result is bounded by the xi\_ion x clumping x proxy-calibration systematic, rather than statistical errors. Further research is needed to refine these estimates and better understand the underlying physical processes driving reionization. Additionally, incorporating data from upcoming surveys and improving calibration techniques will be essential for reducing systematics and obtaining more accurate results.
\section*{References}
\footnotesize
[Muoz2024] Muñoz et al. (2024). Reionization after JWST: a photon budget crisis?. bibcode 2024MNRAS.535L..37M \par [Park2022] Park et al. (2022). Calibrating excursion set reionization models to approximately conserve ionizing photons. bibcode 2022MNRAS.517..192P \par [Duncan2015] Duncan et al. (2015). Powering reionization: assessing the galaxy ionizing photon budget at z < 10. bibcode 2015MNRAS.451.2030D \par [Davies2021] Davies et al. (2021). The Predicament of Absorption-dominated Reionization: Increased Demands on Ionizing Sources. bibcode 2021ApJ...918L..35D \par [Madau2017] Madau (2017). Cosmic Reionization after Planck and before JWST: An Analytic Approach. bibcode 2017ApJ...851...50M

\end{document}
